% Chapter 0: 引言 - 文献概述
% Bayesian Data Analysis - Gelman et al.

\chapter*{引言：文献概述}
\addcontentsline{toc}{chapter}{引言：文献概述}
\label{chap:introduction}

\section{什么是贝叶斯数据分析}

贝叶斯数据分析是一种统计推断方法，其核心思想是使用概率来量化不确定性。在贝叶斯框架下，所有未知量——包括参数、潜在变量、甚至未来观测——都被视为随机变量，用概率分布来描述。

与传统频率主义统计不同，贝叶斯方法明确使用概率来表达我们对未知量的信念（belief）。这种方法的主要优势包括：

\begin{itemize}
  \item \textbf{自然的不确定性量化}：贝叶斯概率区间可以直接解释为"未知量包含在该区间内的概率"
  \item \textbf{灵活的建模能力}：可以构建复杂的多层次模型
  \item \textbf{原则性的推断框架}：贝叶斯定理为所有推断提供了统一的基础
  \item \textbf{先验信息的整合}：允许将领域知识融入统计分析
\end{itemize}

贝叶斯数据分析的过程可以概括为三个步骤：
\begin{enumerate}
  \item 建立完整的概率模型
  \item 基于观测数据计算后验分布
  \item 评估模型的拟合度和敏感性
\end{enumerate}

\section{核心概念}

\subsection{先验分布（Prior Distribution）}

先验分布代表在观测数据之前，我们对未知参数 $\theta$ 的已有知识或信念。先验分布可以是：
\begin{itemize}
  \item \textbf{信息先验}：基于领域知识或历史数据
  \item \textbf{非信息先验}：尽量减少先验对结果的影响（如 Jeffreys 先验）
  \item \textbf{弱信息先验}：包含一定信息但不完全反映领域知识
\end{itemize}

\subsection{后验分布（Posterior Distribution）}

后验分布是贝叶斯推断的核心——在观测到数据 $y$ 后，关于参数 $\theta$ 的更新后的知识：
\[
p(\theta|y) = \frac{p(y|\theta) \cdot p(\theta)}{p(y)}
\]
其中 $p(y|\theta)$ 是似然函数，$p(y)$ 是边际似然（归一化常数）。

\subsection{预测分布（Predictive Distribution）}

贝叶斯方法可以自然地进行预测推断：
\begin{itemize}
  \item \textbf{先验预测分布}：$p(y) = \int p(y|\theta) p(\theta) d\theta$
  \item \textbf{后验预测分布}：$p(\tilde{y}|y) = \int p(\tilde{y}|\theta) p(\theta|y) d\theta$
\end{itemize}

\section{文献目录}\label{sec:literature}

\subsection{核心教材}

\begin{enumerate}
  \item \textbf{\textit{Bayesian Data Analysis, Third Edition}} — Gelman, Carlin, Stern, Dunson, Vehtari, Rubin

  本课程的核心教材，656页，系统介绍了贝叶斯数据分析的理论与实践。

  \begin{itemize}
    \item Part I: 贝叶斯推断基础（概率、三步骤、单/多参数模型、分层模型）
    \item Part II: 数据分析基础（模型检查、评估、决策分析）
    \item Part III: 高级计算方法（MCMC、变分推断）
    \item Part IV: 回归模型
    \item Part V: 非参数模型
  \end{itemize}
\end{enumerate}

\subsection{课件}

\begin{enumerate}
  \item \textbf{Three Steps of Bayesian Data Analysis} — Chapter 1 课件

  \item \textbf{Single-parameter Models} — Chapter 2 课件

  \item \textbf{Multi-parameter Models} — Chapter 3 课件

  \item \textbf{Hierarchical Models} — Chapter 5 课件

  \item \textbf{Model Checking} — Chapter 6 课件

  \item \textbf{Bayesian Workflow} — Chapter 10 课件

  \item \textbf{Advanced Topics} — Chapter 11 课件
\end{enumerate}

\subsection{补充材料}

\begin{enumerate}
  \item \textbf{Metropolis 算法简要证明}

  MCMC 方法中 Metropolis 算法的理论基础与证明。

  \item \textbf{R 语言实现抽样基础}

  使用 R 语言实现各种抽样方法的实践指南。

  \item \textbf{p 值补充}

  p 值的定义、解释与局限性分析。

  \item \textbf{拒绝抽样补充}

  拒绝抽样方法的详细介绍与实现。
\end{enumerate}

\subsection{参考教材}

\begin{enumerate}
  \item \textbf{Markov Chain Monte Carlo in Practice} — Gilks, Richardson, Spiegelhalter

  MCMC 方法的实践指南，是 BDA 第11-12章的重要补充。

  \item \textbf{Monte Carlo Statistical Methods} — Robert, Casella

  蒙特卡洛统计方法的理论教材，包含更深入的算法分析。

  \item \textbf{Bayesian Statistical Modelling} — Congdon

  贝叶斯统计建模的实用参考，提供了丰富的建模案例。
\end{enumerate}

\section{课程结构}

本课程将按照 BDA 第三版的结构依次学习：

\begin{itemize}
  \item \textbf{Part I（第1-5章）}：贝叶斯推断基础
  \item \textbf{Part II（第6-9章）}：数据分析基础
  \item \textbf{Part III（第10-13章）}：高级计算方法
  \item \textbf{Part IV（第14-18章）}：回归模型
  \item \textbf{Part V（第19-23章）}：非参数模型
\end{itemize}

\section{学习方法}

\begin{Remark}
  建议的学习方法：
  \begin{enumerate}
    \item 先阅读课件把握整体框架
    \item 深入阅读教材对应章节
    \item 完成课后习题巩固理解
    \item 通过编程实践加深概念
  \end{enumerate}
\end{Remark}

\endinput
